%% ===========================================================================
%%  Image-Driven Estimation of Magnetic Configurations in Nanodots
%%  Companion (direct-problem) article to Mendez-Rondon et al., Mater. Des. 2026
%%
%%  Target journal : Materials & Design (Elsevier)
%%  Class          : elsarticle
%%  Build          : latexmk -pdf main.tex      (or: make)
%% ===========================================================================
\documentclass[preprint,5p,times,twocolumn]{elsarticle}

\usepackage[T1]{fontenc}
\usepackage[utf8]{inputenc}
\usepackage{amsmath,amssymb,bm}
\usepackage{graphicx}
\usepackage{booktabs}
\usepackage{multirow}
\usepackage{subcaption}
\usepackage{siunitx}
\usepackage[hidelinks]{hyperref}
\usepackage{xcolor}

\graphicspath{{figures/}}

%% --- Draft-note macro: \todo{...} -------------------------------------------
%% Comment out the coloured definition before submission.
%% NOTE: neither \textcolor nor \textbf may wrap the note. Both grab it as an
%% argument through a SHORT internal (\@textcolor, \text@command), so a blank
%% line inside a note aborts the run with "Paragraph ended before ... was
%% complete". \color and \bfseries are declarations: they take no argument and
%% never scan for \par, so multi-paragraph notes are safe.
%% The math branch exists for notes sitting inside equations, e.g.
%% "\mathcal{H} = \todo{full Hamiltonian}", where \bfseries is illegal.
%% Only the TODO label is bold. Setting a whole multi-paragraph note in bold
%% makes an unreadable block in a narrow two-column measure and wrecks the
%% justification, so the body stays at normal weight.
\newcommand{\todo}[1]{%
  \ifmmode
    \text{\color{red}\textbf{[TODO:} #1\textbf{]}}%
  \else
    {\color{red}\textbf{[TODO:} #1\textbf{]}}%
  \fi}

%% --- Notation ---------------------------------------------------------------
\newcommand{\sz}{s_z}                        % magnetisation component
\newcommand{\bth}{\bm{\theta}}               % Hamiltonian parameter vector
\newcommand{\bthhat}{\hat{\bm{\theta}}}      % recovered parameters
\newcommand{\Jt}[1]{\tilde{J}_{#1}}          % exchange constants
\newcommand{\Kdm}{\tilde{K}_{\mathrm{DM}}}   % DMI
\newcommand{\Kan}[1]{\tilde{K}_{\mathrm{an}#1}}
\newcommand{\Hex}{\tilde{H}_{\mathrm{ex}}}
\newcommand{\Tzero}{T^{(0)}}
\newcommand{\Mdisk}{\mathcal{M}}             % nanodot disk mask
\newcommand{\qpeak}{q_{\mathrm{peak}}}
\newcommand{\Cnn}{C_{nn}}
\newcommand{\zvec}{\mathbf{z}}               % Xception latent
\newcommand{\Ecal}{\mathcal{E}}              % encoder

\journal{Materials \& Design}

\begin{document}

\begin{frontmatter}

\title{Image-Driven Estimation of Magnetic Configurations in Nanodots\\
       \todo{final title}}

%% --- Authors -----------------------------------------------------------------
%% Author list and order confirmed. Affiliations a/b/c are taken verbatim from
%% the predecessor paper (Mendez-Rondon et al., Mater. Des. 2026).
\author[a]{Juan Sebasti\'an M\'endez-Rond\'on}
\author[a]{Carlos Andr\'es Canamejoy-Alegr\'ia}
\author[a]{Andr\'es Marino \'Alvarez-Meza\corref{cor1}}
\ead{amalvarezme@unal.edu.co}
\author[a]{Jorge Iv\'an Montes-Monsalve}
\author[b,c]{Jose Dar\'io Agudelo-Giraldo}
\cortext[cor1]{Corresponding author.}

\address[a]{AI-Lab, Signal Processing and Recognition Group, Universidad Nacional
de Colombia, Sede Manizales, 170003 Manizales, Colombia}
\address[b]{Grupo de Investigaci\'on en F\'isica y Matem\'aticas con \'Enfasis en la
Formaci\'on de Ingenieros, Departamento de F\'isica y Matem\'aticas, Facultad de
Ingenier\'ia, Universidad Aut\'onoma de Manizales, 170002 Manizales, Colombia}
\address[c]{Departamento de F\'isica, Facultad de Ciencias Exactas y Naturales,
Universidad de Caldas, 170002 Manizales, Colombia}

%% --- Highlights: 3-5 bullets, <= 85 characters each (Elsevier limit) --------
\begin{highlights}
\item Conditional DDPM generates relaxed magnetic nanodot textures from Hamiltonian parameters.
\item A latent cosine term closes the generation--regression cycle during training.
\item Latent guidance steers sampling without fine-tuning either network.
\item Three disk observables quantify physical fidelity beyond pixel similarity.
\end{highlights}

\begin{abstract}
\todo{Write last. Mirror the predecessor's arc: (i) physical setting and why the
direct problem matters; (ii) the gap --- generative models reproduce texture but
are never held to the inverse model that must read them back; (iii) what we do
--- a conditional DDPM closed against an Xception regressor through a latent
cosine term, in three variants (frozen encoder, joint fine-tune, guidance-only);
(iv) the headline numbers; (v) what it buys the field.}
\end{abstract}

\begin{keyword}
Magnetic nanodots \sep Denoising diffusion probabilistic models \sep
Inverse problems \sep Cycle consistency \sep Latent guidance \sep Spin textures
\end{keyword}

\end{frontmatter}

%% ===========================================================================
%%  SECTION HIERARCHY MIRRORS THE PREDECESSOR PAPER EXACTLY.
%%  See PREDECESSOR_OUTLINE.md for the extracted original structure.
%%    1. Introduction        2. Related work        3. Materials and methods
%%    4. Experimental set-up 5. Results and discussion  6. Conclusions
%% ===========================================================================

\section{Introduction}
\label{sec:intro}

\todo{Four movements, matching the predecessor's introduction.

(1) Magnetic nanodots and their topological spin textures; the extended
Heisenberg Hamiltonian with symmetric exchange, DMI, magnetocrystalline
anisotropy and the Zeeman field \cite{mendezrondon2026regression}.

(2) The inverse problem and its resolution by regression
\cite{kwon2020hamiltonian,mendezrondon2026regression}: an Xception regressor
recovers $\Kdm$, $\Jt{2}$ and $\Tzero$ with $R^2$ of 0.95, 0.91 and 0.84, while
the higher-order exchange constants stay bounded by the identifiability limits
of the $\sz$ projection.

(3) The direct problem. Atomistic Monte Carlo is the ground truth but is slow;
learned generators are fast. Diffusion models \cite{ho2020ddpm,nichol2021improved}
now dominate conditional image synthesis.

(4) The gap this paper closes: a generator judged only by pixel or perceptual
similarity is never held to the inverse model that has to read its output back.
We make that requirement part of the objective. End with the contribution list.}

\subsection{Contributions}
\label{sec:contributions}
\begin{enumerate}
  \item \todo{A cycle-closing objective --- DDPM denoising loss plus a cosine
        term on the 256-d Xception latent --- restricted to a low-noise
        timestep window.}
  \item \todo{Three integration regimes (frozen encoder, joint fine-tune,
        sampling-time latent guidance with no fine-tuning) and the cost of each.}
  \item \todo{A reduced, defensible set of three disk observables
        ($M$, $\Cnn$, $\qpeak$) and the measurement protocol that makes
        simulated and generated configurations comparable.}
\end{enumerate}

%% ===========================================================================
\section{Related work}
\label{sec:related}

\todo{Three strands, in this order.

\emph{Hamiltonian parameter estimation from magnetic images.}
\cite{kwon2020hamiltonian} established the regression formulation;
\cite{mendezrondon2026regression} added the atomistic nanodot setting and RAM
interpretability. Both are inverse-only.

\emph{Generative models for magnetic textures.} \cite{lee2024twisted} is the
closest prior art and the paper to position against: fully-connected generative
kernels producing magnetic domain images of twisted bilayer CrI$_3$ from
Hamiltonian parameters. State plainly what is different here --- a diffusion
backbone on an atomistic nanodot, and, crucially, a generator that is evaluated
and trained against the inverse model rather than in isolation.

\emph{Guidance and conditioning in diffusion models.} Classifier guidance
\cite{dhariwal2021beat} versus classifier-free guidance \cite{ho2022cfg}; why
the former fits here (a pretrained, frozen inverse model already exists and
retraining it would forfeit the published result).}

%% ===========================================================================
\section{Materials and methods}
\label{sec:methods}

\subsection{Atomistic Hamiltonian framework}
\label{sec:hamiltonian}
\todo{Reproduce the Hamiltonian of \cite{mendezrondon2026regression} with
identical notation so the two papers read as one programme.}

\begin{equation}
  \label{eq:hamiltonian}
  \mathcal{H} = \todo{full Hamiltonian}
\end{equation}

\todo{Table 1 analogue: the eight free parameters with sampling ranges, with
$\Jt{1} = \SI{1.0}{\milli\electronvolt}$ fixed as the energy reference ---
$\Tzero$ [K], $\Jt{2}$, $\Jt{3}$, $\Jt{4}$ [meV], $\Kan{1}$, $\Kan{S}$,
$\Hex$ [meV/atom], $\Kdm$ [meV].}

\subsection{GPU-accelerated image dataset generation}
\label{sec:dataset}
\todo{Metropolis--Hastings Monte Carlo with simulated annealing, GPU-accelerated
through the matrix decomposition of \cite{mendezrondon2026regression}. Cylindrical
geometry, $R_d = 18.25$ MUC, thickness $\breve{L} = 5$ MUC, $\sim$5200 localised
spins. The $\sz$ field of the central layer $z = \lfloor \breve{L}/2 \rfloor$ is
projected onto a $39\times39$ pixel grid. Dataset size for this study,
splits 70/15/15, seed 42, and the separate min--max scalers for the generator and
the regressor.}

\subsection{Conditional diffusion framework}
\label{sec:ddpm}
\todo{Forward and reverse processes \cite{ho2020ddpm}; cosine $\beta$-schedule
\cite{nichol2021improved}; U-Net backbone with parameter conditioning injected
through a time--condition embedding; EMA; min-SNR-$\gamma$ weighting
\cite{hang2023minsnr} with $\gamma = 5$.}

\paragraph{Canvas convention.}
The physical image is $39\times39$, which no power-of-two U-Net accepts without
resampling. We reflect-pad the \emph{right and bottom} edges to a $40\times40$
canvas and crop the top-left $39\times39$ block back out at inference. The crop
is exact --- it returns the original pixels --- so the pipeline introduces no
interpolation anywhere. This matters because the annealing temperature $\Tzero$
is encoded in fine-grained texture that bilinear resampling destroys.

\subsection{Cycle integration via the inverse latent}
\label{sec:integration}

\todo{Opening paragraph: the inverse model of
\cite{mendezrondon2026regression} is an Xception \cite{chollet2017xception}
with the shared head
$\mathrm{BN} \to \mathrm{Drop}(0.4) \to \mathrm{Dense}(256,\mathrm{ReLU})
\to \mathrm{BN} \to \mathrm{Drop}(0.3) \to \mathrm{Dense}(8)$.
Removing the final linear layer exposes
$\zvec = \Ecal(x) \in \mathbb{R}^{256}$. Argue why the latent and not the
regression output: the 8-d output is a lossy, task-saturated projection, and
matching it directly would optimise the generator against the regressor's
residual error rather than against the texture information the regressor uses.}

\subsubsection{Latent cosine objective}
\label{sec:cosine}
Let $\hat{x}_0$ be the denoised estimate at timestep $t$ and $x_0$ the reference
configuration, both cropped to $39\times39$. The training objective is

\begin{equation}
  \label{eq:loss}
  \mathcal{L} \;=\; \mathcal{L}_{\text{denoise}}
    \;+\; \lambda_{\cos}\,
    \mathbb{E}\!\left[\mathbb{1}[t < t_{\max}]
    \left(1 - \cos\!\big(\Ecal(\hat{x}_0),\,
    \mathrm{sg}[\Ecal(x_0)]\big)\right)\right],
\end{equation}

where $\mathrm{sg}[\cdot]$ is the stop-gradient and $\mathbb{1}[t < t_{\max}]$
restricts the term to the low-noise window.

\paragraph{Why the timestep window.}
\todo{The methodological point worth defending. At large $t$, $\hat{x}_0$ is a
clamped reconstruction from near-pure noise: its latent carries no usable signal
and the clamp zeroes the gradient. Applying the cosine term across all $t$ spends
compute on a term that is simultaneously uninformative and non-differentiable.
Report the ablation over $t_{\max}$.}

\subsubsection{Frozen encoder and joint fine-tuning}
\label{sec:variants}
\todo{Variant A freezes $\Ecal$: gradients reach $\hat{x}_0$ but no encoder
weight moves. Variant B lets $\Ecal$ train at a reduced learning rate. State the
two guardrails B requires --- latent-collapse monitoring and per-epoch $R^2$ of
the full regressor --- and why: a jointly trained encoder can minimise the cosine
term by degenerating towards a constant embedding, and fidelity bought by
damaging the inverse model is not a gain.}

\subsubsection{Latent-guided sampling}
\label{sec:guidance}
\todo{Variant C changes no weights in either network. A small
$\bth \mapsto \zvec^\star$ head supplies the target embedding, and sampling
follows classifier guidance \cite{dhariwal2021beat} on the cosine objective:

\begin{equation*}
  \hat{\epsilon} \leftarrow \epsilon_\vartheta(x_t, t, \bth)
    + s\,\sqrt{1-\bar\alpha_t}\,
      \nabla_{x_t}\big(1 - \cos(\Ecal(\hat{x}_0), \zvec^\star)\big).
\end{equation*}

Contrast with classifier-free guidance \cite{ho2022cfg}, which would have
required retraining. Report the held-out cosine of the $\bth \mapsto \zvec^\star$
head --- it upper-bounds what guidance can achieve.}

%% ===========================================================================
\section{Experimental set-up}
\label{sec:setup}

\subsection{Atomistic simulation details}
\label{sec:setup-sim}
\todo{Annealing schedule, Monte Carlo sweeps, convergence criterion, and the
number of parameter points. Reuse the predecessor's settings verbatim so the
datasets are identical.}

\subsection{Deep learning training}
\label{sec:setup-training}
\todo{DDPM: $T = 1000$ steps, cosine schedule, base width 80, conditioning
embedding 128, dropout 0.1, EMA 0.999, AdamW $2.26\times10^{-4}$, weight decay
$4.28\times10^{-5}$, batch 128, 3 warm-up epochs, gradient clipping at 0.5,
min-SNR $\gamma = 5$. Cycle terms: $\lambda_{\cos}$, $t_{\max}$, and for the
joint variant the encoder learning rate $10^{-5}$. Hardware and wall-clock cost.}

\subsection{Evaluation protocol}
\label{sec:setup-eval}

\subsubsection{Physical observables}
\label{sec:physical}
All observables are computed on the $39\times39$ crop, restricted to the nanodot
disk $\Mdisk$ of radius $R_d = 18.25$ MUC centred on the image
($N_{\Mdisk} = 1049$ pixels). \emph{The crop precedes the mask}: applying the
disk to the $40\times40$ canvas would average reflected duplicate pixels on a
mask displaced by half a pixel.

\begin{align}
  M      &= \frac{1}{N_\Mdisk}\sum_{i \in \Mdisk} \sz(i),
  \label{eq:magnetization}\\[2pt]
  \Cnn   &= \big\langle \sz(i)\,\sz(j) \big\rangle_{\langle i,j\rangle \in \Mdisk},
  \label{eq:correlation}\\[2pt]
  \qpeak &= \frac{2\pi}{L}\,
            \arg\max_{r>0} \; \big\langle S(\mathbf{q}) \big\rangle_{|\mathbf{q}|=r},
            \qquad L = 39.
  \label{eq:qpeak}
\end{align}

\todo{Justify the reduction to three. Susceptibility and specific heat require
either an ensemble or a single-configuration proxy whose validity is
regime-dependent, so neither is a uniform comparison axis across the magnetic
phases; the exchange energy density is monotonically tied to $\Cnn$ under the
$\sz$ projection. $M$, $\Cnn$ and $\qpeak$ are well defined on every single
configuration in every phase. $S(\mathbf{q})$ is computed on the fluctuation
field --- disk mean removed, background zeroed --- so $\qpeak$ describes the
texture rather than the aperture.}

\subsubsection{Image and regression metrics}
\label{sec:othermetrics}
\todo{SSIM \cite{wang2004ssim} and masked MSE for image fidelity; $R^2$, MAE and
RMSE in physical units for cycle closure, reported per parameter as in the
predecessor's Table 3. Note the known limitation: a phase-shifted texture is
energetically equivalent but penalised by SSIM, which is why the physical
observables and the latent cosine are reported alongside it.}

%% ===========================================================================
\section{Results and discussion}
\label{sec:results}

\subsection{Generative fidelity of the baseline DDPM}
\label{sec:res-baseline}
\todo{Establishes the reference. Figure: generated vs simulated panels across
the six magnetic phases, in the predecessor's phase-diagram style. Table:
per-parameter cycle $R^2$/MAE.}

\subsection{Effect of the latent cosine term}
\label{sec:res-cosine}
\todo{Frozen-encoder variant against baseline. Sweep $\lambda_{\cos}$ and
$t_{\max}$. Report the three observables, SSIM, and cycle $R^2$.}

\subsection{Cost of joint fine-tuning}
\label{sec:res-joint}
\todo{Variant B. Report generative gain \emph{and} inverse-model degradation on
the same axes, plus the latent-collapse diagnostic. If the inverse model
degrades, say so --- that is the result.}

\subsection{Guidance without fine-tuning}
\label{sec:res-guidance}
\todo{Variant C. Sweep $s \in \{0, 0.5, 1, 2, 4, 8\}$; $s = 0$ must reproduce the
baseline exactly. Show where guidance starts trading texture realism for latent
agreement.}

\subsection{Per-phase behaviour}
\label{sec:res-phases}
\todo{Breakdown across the magnetic phases, mirroring the predecessor's Table 4.
That paper showed degenerate textures dominate the inverse error; test whether
the same phases dominate the generative error, and whether the cycle terms help
precisely there.}

\subsection{Limitations}
\label{sec:limitations}
\todo{Be explicit, as the predecessor's Section 5.5 is: $\sz$ projection only,
single-dot geometry, the $\qpeak$ radial-bin quantisation of
$2\pi/39 \approx 0.16$ rad/site, and the fact that the cosine term is bounded by
the encoder's own fidelity.}

%% ===========================================================================
\section{Conclusions}
\label{sec:conclusions}
\todo{Three to five paragraphs. What was built, what the numbers show, which of
the three integration regimes is the recommendation and under what constraint,
and the next step.}

%% ===========================================================================
\section*{CRediT authorship contribution statement}
\todo{Match the predecessor's CRediT block.}

\section*{Declaration of Generative AI and AI-assisted technologies in the writing process}
\todo{The predecessor carries this declaration; Elsevier requires it. State the
tools used and confirm the authors reviewed and take responsibility for the
content.}

\section*{Declaration of competing interest}
The authors declare that they have no known competing financial interests or
personal relationships that could have appeared to influence the work reported
in this paper.

\section*{Acknowledgements}
This work was supported by Hermes 62642, Universidad Nacional de Colombia, and
by Code 111908, Minciencias (951-2024).

\appendix
\section{Supplementary data}
\label{app:supp}
\todo{Additional per-phase tables and the full guidance-scale sweep.}

\section*{Data availability}
The datasets are publicly available on Kaggle
(\texttt{carloscanamejoy/dataset-spines-united-v2}). Code is available at
\url{https://github.com/canamejoy/Spin-State-GenerationAndInference}.

\bibliographystyle{elsarticle-num}
\bibliography{refs}

\end{document}
